% GENERATED FILE. Edit canonical lessons, then run npm run generate:books.
% canonical-source-hash: fnv1a64:c73959393d2be206
% canonical-lessons: HI-C434-asli, HI-C434-suddh, HI-C434-ragarna, HI-C434-phabna, HI-C434-kasauti

\chapter{Pure Gold and Touchstones}
\label{ch:hi-pure-gold-and-touchstones}
\hlchaptermodality{434}

\begin{chapteropening}
\textbf{By the end of this chapter:} \emph{I can tell real from pure, say a colour suits someone and see gold tested on a touchstone.}

Complete the last lesson of chapter 434: I can tell real from pure, say a colour suits someone and see gold tested on a touchstone.
\end{chapteropening}

\section[\texorpdfstring{\hi{असली} (aslī)}{असली (aslī)}]{\hi{असली} (aslī) — real, genuine}

\label{lesson:HI-C434-asli}



\begin{quote}
Before the new one: say the Hindi for a weekly market, then the Hindi for to hand (to).
\end{quote}

\subsection*{You'll want to know: \hi{असली}}
\textbf{\hi{असली}} — \emph{aslī} — \textquotedblleft{}real, genuine\textquotedblright{}.

\begin{quote}
\hi{यह} \hi{असली} \hi{चाँदी} \hi{है।}
\end{quote}

\begin{quote}
\emph{yah aslī cā̃dī hai.}
\end{quote}

\begin{quote}
This is real silver.
\end{quote}

\begin{cousinweb}
\textbf{\hi{असली}} (\emph{aslī}) is usually traced to Arabic \emph{aṣlī}, \textquotedblleft{}original\textquotedblright{}, from \emph{aṣl}, \textquotedblleft{}a root\textquotedblright{}, through Persian: the thing itself, not a copy.
\end{cousinweb}

\begin{grammarlens}[title={What you've built}]
One more word for getting around.
\end{grammarlens}

\subsection*{Your turn}
\begin{itemize}
\raggedright
  \item \emph{Say it:} \emph{aslī}
  \item \emph{Say it:} \emph{aslī}, once more
  \item \emph{Say it:} \emph{yah aslī cā̃dī hai.} — This is real silver.
  \item \emph{Recall:} say the Hindi for a weekly market, then the Hindi for to hand (to), then say \emph{aslī} again
\end{itemize}

\begin{tcolorbox}[breakable,colback=teal!4,colframe=teal!35!black,title={Before you move on}]
What does \hi{असली} mean? (\textquotedblleft{}Real, genuine\textquotedblright{}.) Say it once more.
\end{tcolorbox}
\section[\texorpdfstring{\hi{शुद्ध} (śuddh)}{शुद्ध (śuddh)}]{\hi{शुद्ध} (śuddh) — pure; unmixed}

\label{lesson:HI-C434-suddh}



\begin{quote}
Before the new one: say the Hindi for to hand (to), then the Hindi for real.
\end{quote}

\subsection*{You'll want to know: \hi{शुद्ध}}
\textbf{\hi{शुद्ध}} — \emph{śuddh} — \textquotedblleft{}pure; unmixed\textquotedblright{}.

\begin{quote}
\hi{सुनार} \hi{ने} \hi{कहा} \hi{कि} \hi{यह} \hi{शुद्ध} \hi{सोना} \hi{है।}
\end{quote}

\begin{quote}
\emph{sunār ne kahā ki yah śuddh sonā hai.}
\end{quote}

\begin{quote}
The goldsmith said that this is pure gold.
\end{quote}

\begin{cousinweb}
\textbf{\hi{शुद्ध}} (\emph{śuddh}) is a tatsama, Sanskrit \emph{śuddha}, \textquotedblleft{}cleansed, clear, pure\textquotedblright{} (Monier-Williams): gold with nothing mixed in, or water fit to drink.
\end{cousinweb}

\begin{grammarlens}[title={What you've built}]
One more word for getting around.
\end{grammarlens}

\subsection*{Your turn}
\begin{itemize}
\raggedright
  \item \emph{Say it:} \emph{śuddh}
  \item \emph{Say it:} \emph{śuddh}, once more
  \item \emph{Say it:} \emph{sunār ne kahā ki yah śuddh sonā hai.} — The goldsmith said that this is pure gold.
  \item \emph{Recall:} say the Hindi for to hand (to), then the Hindi for real, then say \emph{śuddh} again
\end{itemize}

\begin{tcolorbox}[breakable,colback=teal!4,colframe=teal!35!black,title={Before you move on}]
What does \hi{शुद्ध} mean? (\textquotedblleft{}Pure; unmixed\textquotedblright{}.) Say it once more.
\end{tcolorbox}
\section[\texorpdfstring{\hi{रगड़ना} (ragaṛnā)}{रगड़ना (ragaṛnā)}]{\hi{रगड़ना} (ragaṛnā) — to rub}

\label{lesson:HI-C434-ragarna}



\begin{quote}
Before the new one: say the Hindi for real, then the Hindi for pure.
\end{quote}

\subsection*{You'll want to know: \hi{रगड़ना}}
\textbf{\hi{रगड़ना}} — \emph{ragaṛnā} — \textquotedblleft{}to rub\textquotedblright{}.

\begin{quote}
\hi{सुनार} \hi{ने} \hi{सोने} \hi{को} \hi{पत्थर} \hi{पर} \hi{रगड़ा।}
\end{quote}

\begin{quote}
\emph{sunār ne sone ko patthar par ragaṛā.}
\end{quote}

\begin{quote}
The goldsmith rubbed the gold on a stone.
\end{quote}

Past, with \hi{ने} (he or she …): \hi{उसने} \hi{हाथ} \hi{रगड़े}, \emph{usne hāth ragaṛe}.

\begin{cousinweb}
\textbf{\hi{रगड़ना}} (\emph{ragaṛnā}) goes back to \emph{*ragg}, \textquotedblleft{}rub\textquotedblright{}, a form CDIAL reconstructs (CDIAL 10558), which lists Hindi \emph{ragaṛnā}. Goldsmiths test gold by the streak it leaves when rubbed on a dark stone.
\end{cousinweb}

\begin{grammarlens}[title={What you've built}]
One more word for getting around.
\end{grammarlens}

\subsection*{Your turn}
\begin{itemize}
\raggedright
  \item \emph{Say it:} \emph{ragaṛnā}
  \item \emph{Say it:} \emph{ragaṛnā}, once more
  \item \emph{Say it:} \emph{sunār ne sone ko patthar par ragaṛā.} — The goldsmith rubbed the gold on a stone.
  \item \emph{Recall:} say the Hindi for real, then the Hindi for pure, then say \emph{ragaṛnā} again
\end{itemize}

\begin{tcolorbox}[breakable,colback=teal!4,colframe=teal!35!black,title={Before you move on}]
What does \hi{रगड़ना} mean? (\textquotedblleft{}To rub\textquotedblright{}.) Say it once more.
\end{tcolorbox}
\section[\texorpdfstring{\hi{फबना} (phabnā)}{फबना (phabnā)}]{\hi{फबना} (phabnā) — to suit; to look good on}

\label{lesson:HI-C434-phabna}



\begin{quote}
Before the new one: say the Hindi for pure, then the Hindi for to rub.
\end{quote}

\subsection*{You'll want to know: \hi{फबना}}
\textbf{\hi{फबना}} — \emph{phabnā} — \textquotedblleft{}to suit; to look good on\textquotedblright{}.

\begin{quote}
\hi{यह} \hi{लाल} \hi{कुर्ता} \hi{तुम} \hi{पर} \hi{बहुत} \hi{फबता} \hi{है।}
\end{quote}

\begin{quote}
\emph{yah lāl kurtā tum par bahut phabtā hai.}
\end{quote}

\begin{quote}
This red kurta really suits you.
\end{quote}

Past (it …): \hi{वह} \hi{साड़ी} \hi{उस} \hi{पर} \hi{फबी}, \emph{vah sāṛī us par phabī}.

\begin{cousinweb}
\textbf{\hi{फबना}} (\emph{phabnā}) is filed by CDIAL under \emph{*sparvati}, \textquotedblleft{}releases, succeeds, is fitting\textquotedblright{}, a form it reconstructs (CDIAL 13808), with Hindi \emph{phabnā}, \textquotedblleft{}to be fitting\textquotedblright{}.
\end{cousinweb}

\begin{grammarlens}[title={What you've built}]
One more word for getting around.
\end{grammarlens}

\subsection*{Your turn}
\begin{itemize}
\raggedright
  \item \emph{Say it:} \emph{phabnā}
  \item \emph{Say it:} \emph{phabnā}, once more
  \item \emph{Say it:} \emph{yah lāl kurtā tum par bahut phabtā hai.} — This red kurta really suits you.
  \item \emph{Recall:} say the Hindi for pure, then the Hindi for to rub, then say \emph{phabnā} again
\end{itemize}

\begin{tcolorbox}[breakable,colback=teal!4,colframe=teal!35!black,title={Before you move on}]
What does \hi{फबना} mean? (\textquotedblleft{}To suit; to look good on\textquotedblright{}.) Say it once more.
\end{tcolorbox}
\section[\texorpdfstring{\hi{कसौटी} (kasauṭī)}{कसौटी (kasauṭī)}]{\hi{कसौटी} (kasauṭī) — a touchstone; a test}

\label{lesson:HI-C434-kasauti}



\begin{quote}
Before the new one: say the Hindi for to rub, then the Hindi for to suit.
\end{quote}

\subsection*{You'll want to know: \hi{कसौटी}}
\textbf{\hi{कसौटी}} — \emph{kasauṭī} — \textquotedblleft{}a touchstone; a test\textquotedblright{}.

\begin{quote}
\hi{सुनार} \hi{कसौटी} \hi{पर} \hi{सोना} \hi{रगड़कर} \hi{देखता} \hi{है।}
\end{quote}

\begin{quote}
\emph{sunār kasauṭī par sonā ragaṛkar dekhtā hai.}
\end{quote}

\begin{quote}
The goldsmith tests gold by rubbing it on a touchstone.
\end{quote}

Feminine.

\begin{cousinweb}
\textbf{\hi{कसौटी}} (\emph{kasauṭī}) is a tadbhava of Sanskrit \emph{kaṣapaṭṭikā}, \textquotedblleft{}a touchstone\textquotedblright{} (CDIAL 2973), which lists Hindi \emph{kasauṭī}; CDIAL suspects the Sanskrit spelling is itself a learned remaking of an older form. The goldsmith's dark stone also stands for any real test: \hi{समय} \hi{की} \hi{कसौटी}, the test of time.
\end{cousinweb}

\begin{grammarlens}[title={What you've built}]
One more word for getting around.
\end{grammarlens}

\subsection*{Your turn}
\begin{itemize}
\raggedright
  \item \emph{Say it:} \emph{kasauṭī}
  \item \emph{Say it:} \emph{kasauṭī}, once more
  \item \emph{Say it:} \emph{sunār kasauṭī par sonā ragaṛkar dekhtā hai.} — The goldsmith tests gold by rubbing it on a touchstone.
  \item \emph{Recall:} say the Hindi for to rub, then the Hindi for to suit, then say \emph{kasauṭī} again
\end{itemize}

\begin{tcolorbox}[breakable,colback=teal!4,colframe=teal!35!black,title={Before you move on}]
What does \hi{कसौटी} mean? (\textquotedblleft{}A touchstone; a test\textquotedblright{}.) Say it once more.
\end{tcolorbox}
